\documentclass[12pt,a4paper]{article}
\usepackage[utf8]{inputenc}
\usepackage{geometry}
\usepackage{hyperref}
\usepackage{listings}
\usepackage{xcolor}
\usepackage{fontawesome5}

\geometry{margin=1in}
\hypersetup{colorlinks=true, linkcolor=blue, urlcolor=blue}

\lstset{
  basicstyle=\ttfamily\footnotesize,
  keywordstyle=\color{blue},
  commentstyle=\color{gray},
  stringstyle=\color{orange},
  breaklines=true,
  showstringspaces=false
}

\title{Wide World Importers SQL Exercises \\ \large Documentation and Solutions}
\author{ {\large by \href{https://github.com/konradcinkusz}{\faGithub\;dev\_insight}}}


\begin{document}
\maketitle

\tableofcontents
\newpage

\section*{Introduction}
This document collects twenty SQL exercises based on the \textbf{Wide World Importers (WWI)} sample database.
WWI models a fictitious wholesale import business and contains tables such as \texttt{Sales.Customers}, \texttt{Sales.Orders},
\texttt{Warehouse.StockItems}, \texttt{Purchasing.Suppliers}, and \texttt{Application.People}.

The goal of these exercises is to practice real-world SQL queries:
\begin{itemize}
    \item Joins across tables
    \item Aggregation with \texttt{GROUP BY}
    \item Filtering with \texttt{WHERE}/\texttt{HAVING}
    \item Subqueries
    \item Window functions (\texttt{RANK}, \texttt{ROW\_NUMBER})
    \item Set operations such as \texttt{UNION} and \texttt{INTERSECT}
  \end{itemize}

  \section*{Prerequisites}
  To try these queries you need access to the Wide World Importers sample database in a SQL Server environment:
  \begin{itemize}
      \item Install SQL Server (Developer or Express) and SQL Server Management Studio, then restore the WWI backup.
      \item Run Microsoft's SQL Server Docker image and restore the backup inside the container.
      \item Use Azure SQL Database and import the WWI BACPAC.
  \end{itemize}
  \noindent Sample backups and BACPAC files are available from Microsoft:\
  \href{https://learn.microsoft.com/en-us/sql/samples/wide-world-importers-oltp-install-configure}{Install and configure Wide World Importers}.

  \section{Exercise 1 (Beginner): Customer Names and Categories}
\textbf{Task:} List each customer’s name along with the name of their customer category.

\textbf{Explanation:}
\begin{enumerate}
  \item \texttt{Sales.Customers} contains \texttt{CustomerName} and a foreign key \texttt{CustomerCategoryID}.
  \item \texttt{Sales.CustomerCategories} defines \texttt{CustomerCategoryName}.
  \item Join the tables on \texttt{CustomerCategoryID}. Use a \texttt{LEFT JOIN} so customers without a category still appear.
\end{enumerate}

\textbf{SQL:}
\begin{lstlisting}[language=SQL]
SELECT c.CustomerName, cc.CustomerCategoryName
FROM Sales.Customers AS c
LEFT JOIN Sales.CustomerCategories AS cc
  ON c.CustomerCategoryID = cc.CustomerCategoryID;
\end{lstlisting}

\section{Exercise 2 (Beginner): Salespeople in the People Table}
\textbf{Task:} Retrieve the full names and email addresses of all people who are employees and are designated as salespeople.

\textbf{Explanation:}
\begin{enumerate}
  \item The table \texttt{Application.People} stores \texttt{FullName}, \texttt{EmailAddress}, \texttt{IsEmployee}, and \texttt{IsSalesperson}.
  \item Filter for rows where both \texttt{IsEmployee} = 1 and \texttt{IsSalesperson} = 1.
\end{enumerate}

\textbf{SQL:}
\begin{lstlisting}[language=SQL]
SELECT FullName, EmailAddress
FROM Application.People
WHERE IsEmployee = 1
  AND IsSalesperson = 1;
\end{lstlisting}

\section{Exercise 3 (Beginner): Supplier Names and Categories}
\textbf{Task:} List each supplier’s name along with the name of the supplier category they belong to.

\textbf{Explanation:}
\begin{enumerate}
  \item \texttt{Purchasing.Suppliers} provides \texttt{SupplierName} and a foreign key \texttt{SupplierCategoryID}.
  \item \texttt{Purchasing.SupplierCategories} contains \texttt{SupplierCategoryName}.
  \item Join the tables on \texttt{SupplierCategoryID} to pair suppliers with category names.
\end{enumerate}

\textbf{SQL:}
\begin{lstlisting}[language=SQL]
SELECT s.SupplierName, sc.SupplierCategoryName
FROM Purchasing.Suppliers AS s
LEFT JOIN Purchasing.SupplierCategories AS sc
  ON s.SupplierCategoryID = sc.SupplierCategoryID;
\end{lstlisting}

\section{Exercise 4 (Beginner): Order Count per Customer}
\textbf{Task:} Find the total number of orders placed by each customer. Return the customer name and the count of orders.

\textbf{Explanation:}
\begin{enumerate}
  \item \texttt{Sales.Orders} records \texttt{OrderID} and \texttt{CustomerID}.
  \item Join \texttt{Sales.Orders} to \texttt{Sales.Customers} to retrieve \texttt{CustomerName}.
  \item Group by \texttt{CustomerName} and count \texttt{OrderID} for each customer.
  \item Optionally order the result by the count to highlight the most active customers.
\end{enumerate}

\textbf{SQL:}
\begin{lstlisting}[language=SQL]
SELECT c.CustomerName,
       COUNT(o.OrderID) AS TotalOrders
FROM Sales.Customers AS c
LEFT JOIN Sales.Orders AS o
  ON c.CustomerID = o.CustomerID
GROUP BY c.CustomerName
ORDER BY TotalOrders DESC;
\end{lstlisting}

\section{Exercise 5 (Beginner): Customers with More Than 5 Orders}
\textbf{Task:} List the names of customers who have placed more than five orders.

\textbf{Explanation:}
\begin{enumerate}
  \item Join \texttt{Sales.Orders} to \texttt{Sales.Customers} on \texttt{CustomerID}.
  \item Group by \texttt{CustomerName} to count orders per customer.
  \item Use \texttt{HAVING COUNT(OrderID) > 5} to keep only customers with more than five orders.
\end{enumerate}

\textbf{SQL:}
\begin{lstlisting}[language=SQL]
SELECT c.CustomerName,
       COUNT(o.OrderID) AS OrderCount
FROM Sales.Customers AS c
JOIN Sales.Orders AS o
  ON c.CustomerID = o.CustomerID
GROUP BY c.CustomerName
HAVING COUNT(o.OrderID) > 5;
\end{lstlisting}

\section{Exercise 6 (Intermediate): Customers Ordering Large Quantities}
\textbf{Task:} Find the names of customers who have ordered more than 100 units of any single stock item in an order line.

\textbf{Explanation:}
\begin{enumerate}
  \item Link \texttt{Sales.Customers} to \texttt{Sales.Orders} and then to \texttt{Sales.OrderLines}.
  \item \texttt{Sales.OrderLines} contains the \texttt{Quantity} of each item in an order.
  \item Filter for \texttt{Quantity > 100} to identify large purchases.
  \item Select distinct customer names so each qualifying customer appears once.
\end{enumerate}

\textbf{SQL:}
\begin{lstlisting}[language=SQL]
SELECT DISTINCT c.CustomerName
FROM Sales.Customers AS c
JOIN Sales.Orders AS o
  ON c.CustomerID = o.CustomerID
JOIN Sales.OrderLines AS ol
  ON o.OrderID = ol.OrderID
WHERE ol.Quantity > 100;
\end{lstlisting}

\section{Exercise 7 (Intermediate): Top 3 Best-Selling Stock Items by Quantity}
\textbf{Task:} Using a window function, list the top three stock items that have the highest total quantity ordered across all sales. Include the total quantity and a rank.

\textbf{Explanation:}
\begin{enumerate}
  \item Aggregate \texttt{Sales.OrderLines} by \texttt{StockItemID} and sum \texttt{Quantity}.
  \item Apply \texttt{RANK() OVER (ORDER BY SUM(Quantity) DESC)} to assign ranks based on quantity sold.
  \item Join to \texttt{Warehouse.StockItems} to obtain \texttt{StockItemName}.
  \item Filter for rows where the rank is less than or equal to three.
\end{enumerate}

\textbf{SQL:}
\begin{lstlisting}[language=SQL]
WITH ItemTotals AS (
  SELECT 
    ol.StockItemID,
    SUM(ol.Quantity) AS TotalQty,
    RANK() OVER (ORDER BY SUM(ol.Quantity) DESC) AS QtyRank
  FROM Sales.OrderLines AS ol
  GROUP BY ol.StockItemID
)
SELECT 
  si.StockItemName,
  it.TotalQty,
  it.QtyRank
FROM ItemTotals AS it
JOIN Warehouse.StockItems AS si
  ON it.StockItemID = si.StockItemID
WHERE it.QtyRank <= 3
ORDER BY it.QtyRank;
\end{lstlisting}

\section{Exercise 8 (Intermediate): Customers with No Orders}
\textbf{Task:} List the names of customers who have never placed an order.

\textbf{Explanation:}
\begin{enumerate}
  \item Retrieve all customers from \texttt{Sales.Customers}.
  \item Gather all \texttt{CustomerID} values appearing in \texttt{Sales.Orders}.
  \item Select customers whose ID does not appear in the orders list using \texttt{NOT IN}.
\end{enumerate}

\textbf{SQL:}
\begin{lstlisting}[language=SQL]
SELECT CustomerName
FROM Sales.Customers
WHERE CustomerID NOT IN (
  SELECT CustomerID FROM Sales.Orders
);
\end{lstlisting}

\section{Exercise 9 (Intermediate): Five Most Recent Orders}
\textbf{Task:} Retrieve the five most recent sales orders and show the order ID, customer ID, and order date.

\textbf{Explanation:}
\begin{enumerate}
  \item Use \texttt{Sales.Orders} which stores \texttt{OrderID}, \texttt{CustomerID}, and \texttt{OrderDate}.
  \item Sort the rows by \texttt{OrderDate} in descending order so the newest orders appear first.
  \item Take only the first five rows.
\end{enumerate}

\textbf{SQL:}
\begin{lstlisting}[language=SQL]
SELECT OrderID, CustomerID, OrderDate
FROM Sales.Orders
ORDER BY OrderDate DESC
FETCH FIRST 5 ROWS ONLY;
\end{lstlisting}

\section{Exercise 10 (Intermediate): Total Invoiced Amount per Customer in 2013}
\textbf{Task:} Calculate the total extended price of all invoices for each customer for invoices dated in 2013. Return the customer name and total amount.

\textbf{Explanation:}
\begin{enumerate}
  \item \texttt{Sales.InvoiceLines} includes \texttt{InvoiceID} and \texttt{ExtendedPrice}.
  \item Join to \texttt{Sales.Invoices} to access \texttt{CustomerID} and \texttt{InvoiceDate}.
  \item Join to \texttt{Sales.Customers} for the \texttt{CustomerName}.
  \item Restrict \texttt{InvoiceDate} to the year 2013.
  \item Group by \texttt{CustomerName} and sum \texttt{ExtendedPrice}.
\end{enumerate}

\textbf{SQL:}
\begin{lstlisting}[language=SQL]
SELECT 
  c.CustomerName,
  SUM(il.ExtendedPrice) AS TotalAmount2013
FROM Sales.InvoiceLines AS il
JOIN Sales.Invoices AS i
  ON il.InvoiceID = i.InvoiceID
JOIN Sales.Customers AS c
  ON i.CustomerID = c.CustomerID
WHERE i.InvoiceDate >= '2013-01-01'
  AND i.InvoiceDate < '2014-01-01'
GROUP BY c.CustomerName
ORDER BY TotalAmount2013 DESC;
\end{lstlisting}

\section{Exercise 11 (Beginner): Customers and Buying Groups}
\textbf{Task:} List each customer with the name of the buying group they belong to, if any.

\textbf{Explanation:}
\begin{enumerate}
  \item \texttt{Sales.Customers} holds \texttt{CustomerName} and a nullable \texttt{BuyingGroupID}.
  \item \texttt{Sales.BuyingGroups} stores \texttt{BuyingGroupName}.
  \item Use a \texttt{LEFT JOIN} on \texttt{BuyingGroupID} so customers without a buying group still appear.
\end{enumerate}

\textbf{SQL:}
\begin{lstlisting}[language=SQL]
SELECT C.CustomerName, B.BuyingGroupName
FROM Sales.Customers AS C
LEFT JOIN Sales.BuyingGroups AS B
  ON C.BuyingGroupID = B.BuyingGroupID;
\end{lstlisting}

\section{Exercise 12 (Beginner): Orders Placed in 2013}
\textbf{Task:} Show every order ID and order date from 2013 together with the customer name.

\textbf{Explanation:}
\begin{enumerate}
  \item The table \texttt{Sales.Orders} contains \texttt{OrderID}, \texttt{OrderDate}, and \texttt{CustomerID}.
  \item Join to \texttt{Sales.Customers} to get \texttt{CustomerName}.
  \item Filter orders where the year of \texttt{OrderDate} equals 2013.
\end{enumerate}

\textbf{SQL:}
\begin{lstlisting}[language=SQL]
SELECT O.OrderID, O.OrderDate, C.CustomerName
FROM Sales.Orders AS O
JOIN Sales.Customers AS C
  ON O.CustomerID = C.CustomerID
WHERE YEAR(O.OrderDate) = 2013;
\end{lstlisting}

\section{Exercise 13 (Intermediate): Orders Per Customer Over Five}
\textbf{Task:} Determine how many orders each customer has placed and return only those with more than five orders.

\textbf{Explanation:}
\begin{enumerate}
  \item Join \texttt{Sales.Orders} to \texttt{Sales.Customers}.
  \item Group by \texttt{CustomerName} to count orders per customer.
  \item Use \texttt{HAVING COUNT(*) > 5} to filter after grouping.
\end{enumerate}

\textbf{SQL:}
\begin{lstlisting}[language=SQL]
SELECT C.CustomerName, COUNT(*) AS OrderCount
FROM Sales.Orders AS O
JOIN Sales.Customers AS C
  ON O.CustomerID = C.CustomerID
GROUP BY C.CustomerName
HAVING COUNT(*) > 5;
\end{lstlisting}

\section{Exercise 14 (Intermediate): Customer Categories with More Than 50 Customers}
\textbf{Task:} List each customer category and the number of customers in it, but only include categories with more than 50 customers.

\textbf{Explanation:}
\begin{enumerate}
  \item \texttt{Sales.Customers} includes \texttt{CustomerCategoryID}.
  \item Join to \texttt{Sales.CustomerCategories} for \texttt{CustomerCategoryName}.
  \item Group by \texttt{CustomerCategoryName} and count customers in each group.
  \item Filter with \texttt{HAVING COUNT(*) > 50}.
\end{enumerate}

\textbf{SQL:}
\begin{lstlisting}[language=SQL]
SELECT CC.CustomerCategoryName,
       COUNT(*) AS NumCustomers
FROM Sales.Customers AS C
JOIN Sales.CustomerCategories AS CC
  ON C.CustomerCategoryID = CC.CustomerCategoryID
GROUP BY CC.CustomerCategoryName
HAVING COUNT(*) > 50;
\end{lstlisting}

\section{Exercise 15 (Advanced): Rank Stock Items by Total Revenue}
\textbf{Task:} Rank stock items by total sales revenue and show the top three items with their revenue.

\textbf{Explanation:}
\begin{enumerate}
  \item \texttt{Sales.InvoiceLines} records \texttt{Quantity} and \texttt{UnitPrice} for each sold item.
  \item Join to \texttt{Warehouse.StockItems} to obtain item names.
  \item Group by stock item and compute \texttt{SUM(Quantity * UnitPrice)} as \texttt{TotalRevenue}.
  \item Use \texttt{RANK() OVER (ORDER BY TotalRevenue DESC)} to assign revenue ranks.
  \item Select only rows where the rank is less than or equal to three.
\end{enumerate}

\textbf{SQL:}
\begin{lstlisting}[language=SQL]
WITH RankedItems AS (
    SELECT SI.StockItemName,
           RANK() OVER (ORDER BY SUM(IL.Quantity * IL.UnitPrice) DESC) AS RevRank,
           SUM(IL.Quantity * IL.UnitPrice) AS TotalRevenue
    FROM Sales.InvoiceLines AS IL
    JOIN Warehouse.StockItems AS SI
      ON IL.StockItemID = SI.StockItemID
    GROUP BY SI.StockItemName
)
SELECT StockItemName, TotalRevenue
FROM RankedItems
WHERE RevRank <= 3;
\end{lstlisting}

\section{Exercise 16 (Advanced): Most Recent Order per Customer}
\textbf{Task:} For each customer, find the date of their most recent order.

\textbf{Explanation:}
\begin{enumerate}
  \item Query \texttt{Sales.Customers} as the outer query.
  \item Use a correlated subquery on \texttt{Sales.Orders} to return \texttt{MAX(OrderDate)} for each customer.
  \item The subquery references the outer query’s \texttt{CustomerID}.
\end{enumerate}

\textbf{SQL:}
\begin{lstlisting}[language=SQL]
SELECT C.CustomerName,
       (SELECT MAX(O.OrderDate)
        FROM Sales.Orders AS O
        WHERE O.CustomerID = C.CustomerID
       ) AS LastOrderDate
FROM Sales.Customers AS C;
\end{lstlisting}

\section{Exercise 17 (Advanced): Duplicate Email Addresses}
\textbf{Task:} Identify any duplicate email addresses in the \texttt{Application.People} table.

\textbf{Explanation:}
\begin{enumerate}
  \item Group \texttt{Application.People} by \texttt{EmailAddress} and keep those with a count greater than one.
  \item Join this list of duplicate emails back to \texttt{Application.People} to display the affected rows.
\end{enumerate}

\textbf{SQL:}
\begin{lstlisting}[language=SQL]
SELECT P.PersonID, P.FullName, P.EmailAddress
FROM Application.People AS P
JOIN (
    SELECT EmailAddress
    FROM Application.People
    GROUP BY EmailAddress
    HAVING COUNT(*) > 1
) AS dup
  ON P.EmailAddress = dup.EmailAddress;
\end{lstlisting}

\section{Exercise 18 (Advanced): All Cities Used by Customers or Suppliers}
\textbf{Task:} Retrieve all distinct city names that appear as delivery cities for either customers or suppliers.

\textbf{Explanation:}
\begin{enumerate}
  \item Select \texttt{DeliveryCityID} from \texttt{Sales.Customers}.
  \item Select \texttt{DeliveryCityID} from \texttt{Purchasing.Suppliers}.
  \item Combine the city IDs using \texttt{UNION}.
  \item Return the corresponding \texttt{CityName} values from \texttt{Application.Cities}.
\end{enumerate}

\textbf{SQL:}
\begin{lstlisting}[language=SQL]
SELECT DISTINCT C.CityName
FROM Application.Cities AS C
WHERE C.CityID IN (
    SELECT DeliveryCityID FROM Sales.Customers
    UNION
    SELECT DeliveryCityID FROM Purchasing.Suppliers
);
\end{lstlisting}

\section{Exercise 19 (Advanced): Cities Common to Customers and Suppliers}
\textbf{Task:} List city names that appear in both customer and supplier delivery addresses.

\textbf{Explanation:}
\begin{enumerate}
  \item Select \texttt{DeliveryCityID} from \texttt{Sales.Customers}.
  \item Select \texttt{DeliveryCityID} from \texttt{Purchasing.Suppliers}.
  \item Use \texttt{INTERSECT} to find city IDs present in both sets.
  \item Retrieve the \texttt{CityName} for those IDs from \texttt{Application.Cities}.
\end{enumerate}

\textbf{SQL:}
\begin{lstlisting}[language=SQL]
SELECT CityName
FROM Application.Cities
WHERE CityID IN (
    SELECT DeliveryCityID FROM Sales.Customers
    INTERSECT
    SELECT DeliveryCityID FROM Purchasing.Suppliers
);
\end{lstlisting}

\section{Exercise 20 (Advanced): Stock Items and Groups}
\textbf{Task:} Retrieve every stock item along with its stock group name. Include items with no group and groups with no items.

\textbf{Explanation:}
\begin{enumerate}
  \item \texttt{Warehouse.StockItems} lists the available items.
  \item \texttt{Warehouse.StockItemStockGroups} links items to groups.
  \item \texttt{Warehouse.StockGroups} defines group names.
  \item A \texttt{FULL OUTER JOIN} between items and the link table keeps items without groups.
  \item Another \texttt{FULL OUTER JOIN} to \texttt{Warehouse.StockGroups} keeps groups without items.
\end{enumerate}

\textbf{SQL:}
\begin{lstlisting}[language=SQL]
SELECT SI.StockItemName, SG.StockGroupName
FROM Warehouse.StockItems AS SI
FULL OUTER JOIN Warehouse.StockItemStockGroups AS SIG
  ON SI.StockItemID = SIG.StockItemID
FULL OUTER JOIN Warehouse.StockGroups AS SG
  ON SIG.StockGroupID = SG.StockGroupID;
\end{lstlisting}

\end{document}
